\documentclass[10pt]{article}

\usepackage[margin=1in]{geometry}
\usepackage{amsmath,amssymb,amsthm}
\usepackage{enumitem}
\usepackage{booktabs}
\usepackage[colorlinks=true,linkcolor=blue,citecolor=blue,urlcolor=blue]{hyperref}

\newtheorem{theorem}{Theorem}[section]
\newtheorem{lemma}[theorem]{Lemma}
\newtheorem{corollary}[theorem]{Corollary}
\theoremstyle{definition}
\newtheorem{definition}[theorem]{Definition}
\newtheorem{example}[theorem]{Example}
\theoremstyle{remark}
\newtheorem{remark}[theorem]{Remark}

\newcommand{\R}{\mathbb{R}}
\newcommand{\E}{\mathbb{E}}

\title{Consequentialist Foundations \\[1ex]
{\large Why anything that pursues goals well comes to look like an expected utility maximizer}}
\author{Lecture notes, Agent Foundations module \\ Iliad Intensive}
\date{June 2026}

\begin{document}

\maketitle

\begin{abstract}
These notes are about a cluster of mathematical results, the coherence theorems, that try to say something forced and general about what it looks like to pursue goals effectively. They assume no prior acquaintance with decision theory, utility functions, probability as a theory of belief, or any of the surrounding vocabulary; every term is built from a concrete example before it is named. We start from a practical worry. If powerful goal-pursuing machines are coming, and if we cannot run experiments on them before they exist, then we would like to know in advance what constraints their behavior must satisfy no matter how they are built. The coherence theorems are one answer. They show, in stages, that a system which never works against itself in a certain precise sense can always be described, from the outside, as if it assigned numerical values to outcomes (a ``utility function''), assigned numerical weights to uncertain possibilities obeying the rules of probability, and chose so as to maximize the weighted average value. We develop this in three movements: first the intuitive version, where a system with circular preferences can be drained of resources like a vending machine that gives change for its own change, and avoiding that drain turns out to be exactly a consistency requirement; second a problem the intuitive version hides, namely that ``working against itself'' only makes sense relative to some resource that plays the role of money, and it is not obvious what that resource is for an arbitrary physical object; third the rigorous version, the complete class theorem, which reaches the same conclusion using nothing but the idea that you should not choose an option when a different available option is at least as good in every possible situation and better in some. We end by saying carefully what these theorems do and do not establish, because they are routinely over-read: they describe the shadow that effective goal-pursuit casts on behavior, and say nothing directly about the machinery casting it.
\end{abstract}

\tableofcontents

\section{The problem: saying something about a mind before it exists}
\label{sec:intro}

Here is a situation that is easy to state and hard to be in. Suppose that, sometime in the future, people succeed in building a machine that is very good at bringing about outcomes it is built to bring about, and suppose it is better at this than we are. We would like, before switching it on, to be able to say something trustworthy about how it will behave, because once it is running and out in the world it may be too late to change our minds. The ordinary way we learn how a complicated thing behaves is to build it, watch it, and fix what goes wrong. That method is not available here, both because the thing does not exist yet and because ``watch it and fix what goes wrong'' is exactly the loop we may not get to run a second time.

What could we possibly say in advance about a machine we have not built, whose internal workings we may not even have chosen yet? Not much about the specifics. But notice that some statements about a goal-pursuing system do not depend on its specifics at all. ``It will not pour its resources into a course of action that predictably circles back to where it started, gaining nothing.'' ``If it reliably gets what it is after, then from the outside its choices will hang together as though it valued some outcomes more than others, by consistent amounts.'' Claims like these are not about the circuitry. They are claims about what \emph{effective pursuit of goals} has to look like, on pain of the system tripping over its own feet. If we could prove a few claims of that kind, rigorously, we would have a foothold: properties that hold for the system we are worried about \emph{because} it pursues goals well, whatever else is true of it.

This is the entire ambition of the coherence theorems, and it is worth being honest that the ambition is modest in one direction and surprising in another. Modest, because these theorems will not tell us \emph{what} a system wants, only what the structure of effective wanting must be. Surprising, because that structure turns out to be quite specific. We will find that an effective goal-pursuer can always be described as though it carries around two things: a way of scoring outcomes by how much it wants them, and a way of weighting uncertain possibilities by how likely it takes them to be, combined by a particular kind of averaging. The names for these, once we have built them up, are a \emph{utility function} and a \emph{probability distribution}, and the combining rule is called \emph{maximizing expected utility}. None of these names should mean anything to you yet; we will earn each one from an example.

A warning about a word that will do a great deal of work. When these theorems conclude that a system ``behaves like'' an expected utility maximizer, they mean that its outward choices can be \emph{summarized} by one, the way the motion of a thrown rock can be summarized by ``it minimizes a certain quantity called action'' without the rock containing a tiny accountant computing integrals. The theorems pin down the summary, not the mechanism. This is the source of their strength, since a summary that holds whatever the mechanism is, is a summary we can rely on without knowing the mechanism; and it is the source of their main limitation, which we will come back to at the end.

We proceed in the order that makes the ideas easiest to acquire, which means starting with preferences among things you get for certain (Section~\ref{sec:moneypump}), moving to choices among uncertain prospects and watching probabilities appear (Section~\ref{sec:probabilities}), pausing on the hidden assumption that all of this leans on (Section~\ref{sec:measuringstick}), then giving the clean rigorous theorem (Section~\ref{sec:completeclass}), and finally taking care to state what has and has not been shown (Section~\ref{sec:scope}). The only thing you need to bring is comfort with simple arithmetic and a willingness to take an example seriously before its moral arrives.

\section{Circular wanting, and how it drains you}
\label{sec:moneypump}

Begin with a system that has likes and dislikes but no discipline imposed on them. By ``likes and dislikes'' nothing fancy is meant: the system, offered a choice between two things, tends to take one rather than the other, and we summarize that tendency by saying it \emph{prefers} the one it takes. To keep it concrete, give it opinions about pizza toppings. It prefers onion to pineapple, meaning that offered the two it takes onion. It prefers pineapple to mushroom. And, here is where trouble enters, it prefers mushroom to onion. Each of these three preferences is perfectly ordinary on its own. Together they form a loop: onion is taken over pineapple, pineapple over mushroom, mushroom over onion, and around again.

Watch what a modestly clever vendor can do to a system whose wants run in a loop like this. The system is currently holding a mushroom slice. The vendor offers a trade: give me one penny, and I will swap your mushroom for a pineapple. Since the system prefers pineapple to mushroom and a penny is a trifle, it pays the penny and takes the pineapple. Before it can eat, the vendor offers another swap, this one free: hand back the pineapple and take an onion instead. Since the system prefers onion to pineapple, it accepts, feeling it has come out ahead at no cost. Then a third free swap: trade the onion for a mushroom. Since it prefers mushroom to onion, it accepts again. Now look at where it stands. It is holding a mushroom slice, the same kind of slice it began with, and it is one penny poorer. The vendor can run the loop again, and again, extracting a penny each time the system completes the circle, until its pockets are empty. The system has been turned into a device that pays to be sent in a circle, a sort of pump with money coming out. Steve Omohundro tells the same story without pizza: an agent that prefers being in Berkeley to San Francisco, San Jose to Berkeley, and San Francisco to San Jose can be kept forever in taxis going around the bay, spending its fare to arrive where it could have stayed for free.

There is a name for what the looping system did wrong, and it is the single most important notion in this section, so we state it carefully.

\begin{definition}[Dominated, and a dominated strategy]
\label{def:dominated}
Say that one available course of action is \emph{dominated} by another if the second does at least as well as the first by the system's own lights in every respect that the system cares about, and strictly better in at least one. A system following a dominated course of action is, by its own standard of better and worse, leaving something on the table: there was an alternative right there that it would itself have judged superior.
\end{definition}

The looping system's behavior was dominated by the most boring alternative imaginable, which is to decline every trade and keep the original mushroom slice. That alternative ends with the same slice and keeps the penny, so it is strictly better, and the elaborate cycle of trades is therefore strictly worse. Nothing in the outside world forced the loss. The system inflicted it on itself, and the cause was precisely the loop in its preferences.

Two features of this little argument carry the whole subject, and both are easy to misread, so sit with them.

The first is that the argument is silent about \emph{what} the system ought to want. We did not show that pineapple ``really'' beats mushroom, or that a well-designed system must like onion. The conclusion is conditional and structural: whatever the system's tastes happen to be, if those tastes contain a loop then the system can be drained, and if the system is not to be drainable then its tastes cannot contain a loop. Removing one specific way of working against yourself turns out to be the same thing as imposing a consistency requirement (no loops) on whatever wants are present. Every coherence argument we meet has this shape. None of them tells a system what to pursue; all of them constrain the bookkeeping that surrounds the pursuing.

The second is the direction the reasoning runs, which is opposite to the usual textbook order. We are not assuming the system has some numerical scoring of toppings and deducing that it avoids loops. We are assuming it avoids loops, the drainable kind, and deducing that its preferences must have the structure a numerical scoring would give them. Here is the connection. Suppose you wanted to encode a loopless ranking of the toppings as numbers, assigning each topping a value so that the system prefers whichever topping has the higher number. Onion above pineapple above mushroom would get, say, the numbers $3, 2, 1$. Such a numerical encoding can never contain a loop, because the numbers themselves do not loop: if onion outranks pineapple and pineapple outranks mushroom then onion's number exceeds mushroom's, full stop. So ``can be encoded by an assigned-number ranking'' and ``free of the loops that let you be drained'' are, for a finite list of outcomes, two descriptions of the very same condition. The assigned-number ranking is what we will call a utility function, and we can already see that it is not an extra assumption smuggled in; it is just a convenient way to write down a consistent set of preferences.

\begin{definition}[Utility function]
\label{def:utility}
A \emph{utility function} is an assignment of a real number to each possible outcome, written $U(\text{outcome})$, such that the system prefers one outcome to another exactly when the first has the larger number. The numbers have no meaning beyond their order at this stage: only ``higher means more preferred'' is being claimed. A set of preferences can be captured by some utility function precisely when it is consistent in the loop-free sense above.
\end{definition}

\section{Choosing under uncertainty, and where probabilities come from}
\label{sec:probabilities}

So far every choice has been between outcomes the system receives for sure. Real choices are rarely like that. Usually you choose an \emph{action}, and which outcome the action produces depends on facts you do not control or have not yet worked out: whether it rains, how a coin lands, what some other party does. A choice whose outcome is uncertain in this way is a \emph{gamble}. The remarkable thing, and the heart of this section, is that the same draining argument from the previous section, once it is pushed into the world of gambles, forces a great deal more than mere loop-freeness. It forces the system to behave as though it has attached numbers not only to outcomes but also to the uncertain possibilities themselves, and as though those second numbers obey the familiar rules of probability, and as though the system combines the two kinds of number by a specific averaging. Most strikingly, the two kinds of number have to come into being together; neither can be read off first and the other second.

It helps to state the unusual order of explanation before walking through it. The usual textbook starts by assuming probabilities are already given (somehow), assumes a utility function is already given (somehow), multiplies and adds them, and declares that a rational system maximizes the result. We are going to refuse to start with either. We start only with a system choosing among gambles and the single demand that it not adopt dominated, drainable strategies. From that demand we will extract, as a forced consequence, that the system must be acting as though it multiplies the value of each possible outcome by some numerical weight attached to the possibility of that outcome, and only after seeing what laws those weights are forced to satisfy will we recognize the weights as probabilities and give them the name.

\subsection{Why there must be numerical weights at all}

Take a concrete gamble. Someone offers your system this deal: hand over one apple now, and I will flip a coin; if it lands heads I give you one orange, and if it lands tails I give you three plums. Whether to take the deal cannot be settled by the system's preferences over apples, oranges, and plums alone, because the result is uncertain and the system has to weigh the orange-if-heads branch against the plums-if-tails branch somehow. The claim is that this weighing cannot be arbitrary if the system is to stay undrained.

Suppose, using the utility functions of Definition~\ref{def:utility}, that this system values one orange at a level we will write as $2$ and three plums at a level we will write as $1.5$, in whatever consistent units its utility function uses. (We will write these units with a small euro sign, $\text{\texteuro}2$ and $\text{\texteuro}1.5$, purely to keep them visually distinct from counts of fruit; the euro sign is a bookkeeping mark, not real money.) Suppose further that the system treats the coin as equally likely to fall either way, so that it weights each branch by one half. Then the value it should place on the whole gamble is the half-and-half blend of the two branch values,
\[
\tfrac12 \cdot U(\text{1 orange}) + \tfrac12 \cdot U(\text{3 plums}) \;=\; \tfrac12 \cdot \text{\texteuro}2 + \tfrac12 \cdot \text{\texteuro}1.5 \;=\; \text{\texteuro}1.75 .
\]
A quantity built like this, by taking each possible outcome's value, multiplying it by the weight on that possibility, and adding up, is called an \emph{expected value} of the outcome's utility, or an \emph{expected utility}. The word ``expected'' is a term of art borrowed from gambling and does not mean the system expects to receive exactly $\text{\texteuro}1.75$ of anything; it means this weighted blend is the summary number by which the gamble is to be judged. The claim of the coherence theorems is that a system which is not throwing value away must judge gambles by some such weighted blend, and must choose, among the gambles on offer, the one whose blend comes out highest. A system that scored gambles in any way that could \emph{not} be rewritten as such a weighted blend of outcome values could be handed a menu of gambles on which its choices are dominated, drained exactly as the pizza loop drained the looping system, only with more arithmetic in the proof.

What makes this more than a definition is that the weights are not free. They are forced into a very particular shape.

\subsection{Why the weights have to behave like probabilities}

Call the weights ``probabilities'' provisionally, as a label we have not yet justified, and watch the constraints on them appear one by one out of the no-draining requirement.

The weights across the possible outcomes of a single uncertain event have to add up to one. To see why, suppose a system betting on one flip of a coin puts weight $0.6$ on heads and weight $0.7$ on tails. These add to $1.3$, which looks like nothing worse than enthusiasm but is in fact a leak. Offer the system a contract that pays it $0.8$ apples no matter how the coin lands. By the system's own weights, it scores this contract at $0.6 \cdot \text{\texteuro}0.8 + 0.7 \cdot \text{\texteuro}0.8 = \text{\texteuro}1.04$, more than the $\text{\texteuro}1$ value of a single apple, so it will pay up to a whole apple to obtain a contract that returns it only $0.8$ apples for certain. That is a guaranteed loss of one fifth of an apple, a dominated choice. The only way to make every such guaranteed-loss trade impossible is for the weights across the outcomes of one event, here heads and tails, to add up to exactly one. Whether or not we call them probabilities, the numbers you multiply outcome values by must sum to one.

The packaged form of this kind of argument has a name, the \emph{Dutch book}, and it is worth meeting because it recurs throughout the subject. Picture a bookmaker who must publicly quote, for each event, the price of a ticket that pays one dollar if that event happens. Let $X$ and $Y$ be two events that are mutually exclusive (they cannot both happen) and exhaustive (one of them must), so that exactly one occurs; the bookmaker quotes price $\$x$ for the $X$-ticket and $\$y$ for the $Y$-ticket. Anyone holding both tickets is guaranteed to collect exactly one dollar, since exactly one event occurs. If the two prices do not add to one, a sharp customer can make a sure profit: when $x + y$ exceeds one, sell the pair (you take in more than the dollar you will have to pay out); when it falls short, buy the pair (you pay less than the dollar you are sure to collect). A set of bets rigged to take money from someone no matter what happens is a Dutch book, and the bookmaker avoids being Dutch-booked exactly when the prices of mutually exclusive, exhaustive events sum to one. This is the same fact as the previous paragraph in different dress, and the moral is the same: the no-sure-loss demand forces the weights to add to one.

The structure forced on the weights goes deeper than addition. It pins down how the weight on a combined event relates to the weights on its parts, which is the rule textbooks call conditional probability. Consider two events $Q$ and $R$, and three tickets. The first pays a dollar if $Q$ happens; call its fair price $\$x$. The second is designed to isolate ``$R$, given that $Q$ happened'': it is called off and your money refunded if $Q$ does not happen, and if $Q$ does happen it pays a dollar when $R$ also happens; call its price $\$y$, which is functioning as the price of $R$ \emph{conditional on} $Q$. The third pays a dollar only if $Q$ and $R$ both happen; call its price $\$z$. A customer who finds any inconsistency among these three prices can assemble a guaranteed profit out of them, and working out which price combinations leave no guaranteed profit yields exactly one clean equation,
\[
z \;=\; x \cdot y .
\]
Read $x$ as the weight on $Q$, written $P(Q)$; read $y$ as the weight on $R$ given that $Q$ occurred, written $P(R \mid Q)$ and called the \emph{conditional} weight; read $z$ as the weight on both happening, $P(Q \wedge R)$. The forced equation says $P(Q \wedge R) = P(Q)\,P(R \mid Q)$, which is the standard definition of conditional probability. We did not hand the system this rule; we handed it the no-sure-loss demand and the rule fell out. As a check on how unforgiving the demand is, a bookmaker who quotes $x = 0.60$, $y = 0.70$, and $z = 0.20$ (so that $z \ne x y = 0.42$) can be made to lose for certain by a suitable bundle of the three tickets.

\subsection{What the intuitive picture shows, and what it admits}

Stack the arguments of this section together (the gamble that needs a weighted blend, the weights that must sum to one, the Dutch book, the conditional-probability tickets) and they converge on one description. A system that does as well as it qualitatively can at getting what it is after must be describable as carrying a consistent scoring of outcomes (a utility function), consistent weights on uncertain possibilities that obey the rules of probability (a probability distribution), and a policy of choosing whichever available gamble has the highest probability-weighted score (maximizing expected utility). The probabilities and the utilities are not supplied to the system from outside; they are squeezed out of it, together, by the single demand that it not be drainable. A system that represents its uncertainty with such probability weights and updates them as it learns is what people call a \emph{Bayesian} agent, after Thomas Bayes; the coherence arguments are one route to the claim that effective agents look Bayesian from outside.

It is important to be as candid as Yudkowsky is about the standing of these demonstrations. They are intuition pumps for theorems, not the theorems themselves. They make it vivid that results of this kind exist and roughly why, but they wave their hands at exactly the points a skeptic would press, for instance the precise sense in which a non-blend scoring ``can be drained.'' The rigorous version, which we give in Section~\ref{sec:completeclass}, closes those gaps. Two honest caveats also survive into everything that follows. Perfect coherence is generally impossible to compute, so no real, resource-bounded system is ever exactly coherent; and a system that is a tangled mess on the inside may still be impossible to drain from the outside, so ``looks coherent'' and ``is internally tidy'' are different properties. Both caveats turn out to matter enormously. The first is the entire reason the topic of logical uncertainty exists. The second is the gap between what these theorems say (a fact about outward behavior) and the separate question of what a system is actually built out of.

Before the rigorous theorem, we have to face something every argument in this section quietly relied on each time it said the word ``drained.''

\section{The hidden ruler: what counts as a resource?}
\label{sec:measuringstick}

Every draining argument so far has been keeping its books in some currency. The pizza system lost \emph{pennies}; the gambler lost \emph{apples}; the bookmaker lost \emph{dollars}. ``Strictly worse off'' always meant worse off measured in units of something we had silently agreed was worth accumulating. John Wentworth's observation, which he calls the measuring-stick problem, is that this currency is doing essential work, and that the coherence theorems are quietly leaning on our having one.

The pressure becomes visible through a common objection. People sometimes say the coherence theorems are empty, because \emph{any} system at all can be cast as maximizing some utility function: whatever a rock does, simply declare its utility function to be ``do exactly that,'' and the rock is now a flawless maximizer of rock-utility. If that were the whole story, the theorems would constrain nothing, since everything would satisfy them. The reply is that the trivialization works only if you are allowed to pick the utility function after the fact, fitted to the very behavior you are trying to explain, over the very same variables that behavior ranges over. The coherence theorems are not playing that game. They assume there is some resource in the picture, something that behaves like money, and they pick out the systems that do not squander \emph{that resource}. Relative to a fixed resource, ``dominated strategy'' becomes a real and falsifiable accusation. A system that will pay to be cycled pepperoni to mushroom to anchovy and back to pepperoni is genuinely doing something wrong: either it has no consistent ranking of toppings, or it is violating the assumption that having more money is better. The resource is what gives the theorem teeth instead of letting it collapse into ``everything maximizes something.''

What does a thing have to be like to serve as such a measuring stick? Wentworth isolates two properties, and they are exactly what make an ordinary ruler useful for measuring. First, the resource must add up across separate decisions: paying a penny here and a penny there really costs two pennies, the way two one-centimeter gaps laid end to end really make two centimeters. Second, more of it must be better, all else equal, the way a larger number of centimeters is unambiguously more length. Money has both properties, which is why money is the textbook measuring stick, but the properties are what matter, not the coinage. Underneath them sits a third feature that does the real explanatory work: the resource is \emph{fungible}, meaning it can be traded for many different things the system might want, and it is this tradeability that lets ``how much would you pay'' reveal an entire ranking of preferences. Energy, time, and usable order (the physicists' ``free energy'') are candidate measuring sticks for the same reason, which is one of the threads tying this topic to the optimization-and-thermodynamics material elsewhere in the module: physical processes spend free energy the way shoppers spend money.

The genuinely hard and still-open part is recognizing a measuring stick \emph{in the wild}, in a situation where we do not already think of anything as a resource. When we have already decided that dollars or joules are the currency, the problem is invisible, but only because we imported the answer. Wentworth puts the open question in its sharpest form: handed a simulated world running on some unfamiliar physics, what program could we run on it to pick out, from the physics alone, which quantities are the ``resources''? We do not have such a program. And the prize for having one would be large, because it would let the coherence theorems do work they cannot currently do. If we could read off a system's resources directly from its physics, with no notion of agency assumed in advance, then the coherence theorems would give us a non-circular way to say that some physical systems contain goal-pursuing agents and others do not: a system contains an agent to the extent that it acquires and spends the physically-identified resources the way an undrainable maximizer would. We could point at a region of the world and say, objectively, this much goal-pursuit is happening here, measured against this resource. Wentworth is careful that the intuitive bridge here (effectively steering the world into a narrow set of outcomes almost always requires hoarding broadly useful resources, so systems that fling resources away cannot steer well) is a qualitative argument and not yet a theorem. Turning it into one is open work, and it is one of the places where this topic hands off to the descriptive program described in another note.

The lasting lesson for everything below is a habit of suspicion. Whenever a coherence theorem says ``dominated,'' ask: dominated in what currency, and who chose the currency? The theorem in the next section is admirably explicit about this, because it drops the betting story entirely and lets the currency be whatever an explicitly given scoring of outcomes says it is.

\section{The rigorous version: the complete class theorem}
\label{sec:completeclass}

The arguments of the previous two sections have a feature some people find unsatisfying as a foundation: they define belief as willingness to bet and preference as willingness to pay, and then derive good behavior from the badness of certain bets and payments. The whole justification runs through imagined exploitation, through bookies and money pumps. Abram Demski's presentation of the \emph{complete class theorem} offers a foundation of a different temper, one that compares outcomes directly and never mentions a bet. It reaches the same destination (probabilities, utilities, and the rule of maximizing their weighted blend) from a single starting principle: do not choose a course of action when some other available course is at least as good in every situation you might be in, and strictly better in at least one. We will build the statement from the ground up, defining each ingredient as it appears.

\subsection{The pieces of a decision problem}

To make the principle precise we need a vocabulary for a decision faced under uncertainty. Imagine you must act, you are unsure which of several situations you are actually in, you may first receive a hint, and different actions lead to outcomes of differing goodness in each situation. Five pieces capture this.

\begin{definition}[A decision problem]
\label{def:decprob}
There is a set $\Theta$ of \emph{states of the world}, the different ways reality might actually be, only one of which is true but you do not know which (for instance, ``the patient has the disease'' versus ``the patient does not''). There is a set $X$ of possible \emph{observations}, the hints you might receive before acting (the result of a test). There is a set $A$ of \emph{actions} available to you (treat, or do not treat). There is a \emph{likelihood} $F(x \mid \theta)$, which says how probable the observation $x$ is when the true state is $\theta$ (how often the test reads positive when the patient really is sick). And there is a \emph{loss function} $L(\theta, a)$, a number saying how bad the outcome is if you take action $a$ while the world is in state $\theta$, with smaller numbers meaning better outcomes. A \emph{decision rule} $\delta$ is a complete plan that fixes, for every observation you might see, the action you will then take: $\delta(x)$ is the action the rule prescribes upon seeing $x$.
\end{definition}

Two remarks orient this. The loss function is where the framework keeps its measuring stick from Section~\ref{sec:measuringstick}: it is handed to us as raw material, a primitive scoring of how good each situation-and-action pairing is, and the theorem will manufacture probabilities while taking this scoring as given. (A loss is just a flipped utility, badness instead of goodness, used because the tradition this theorem comes from minimizes loss rather than maximizing gain; nothing turns on the sign.) For the clean result we assume the lists $\Theta$ and $A$ are finite. A decision rule should be judged not by how it fares in one situation but across all the situations it might face, and the right summary of that is its average loss in each state.

\begin{definition}[Risk]
\label{def:risk}
The \emph{risk} of a decision rule $\delta$ in a state $\theta$ is the average loss the rule incurs when the world is in that state, averaging over the observations the state might produce:
\[
R(\theta, \delta) \;=\; \E_{x \sim F(\cdot \mid \theta)}\big[\, L(\theta, \delta(x)) \,\big]
\;=\; \sum_x F(x \mid \theta)\, L\big(\theta, \delta(x)\big).
\]
Collecting one such number for each possible state gives the rule's \emph{risk profile}, a list of numbers $\big(R(\theta, \delta)\big)_{\theta \in \Theta}$, one coordinate per state of the world, which we can picture as a single point in a space with one axis per state.
\end{definition}

\subsection{Not being dominated, made precise}

Now we import, into this cleaner setting, the only normative idea the theorem needs, which is the no-domination idea from Section~\ref{sec:moneypump} translated from ``in every respect I care about'' to ``in every state of the world I might be in.''

\begin{definition}[Pareto improvement and admissibility]
\label{def:admissible}
A rule $\delta^*$ is a \emph{Pareto improvement} over a rule $\delta$ if it has at least as low a risk in every state, $R(\theta, \delta) \ge R(\theta, \delta^*)$ for all $\theta$, and strictly lower risk in at least one state. (The name honors the economist Vilfredo Pareto, who studied changes that help someone and harm no one.) A rule is \emph{admissible} if no Pareto improvement over it exists, and \emph{inadmissible}, or \emph{dominated}, if one does.
\end{definition}

An inadmissible rule is one that some competitor beats in at least one possible situation while never doing worse in any, the across-states analogue of paying a penny to end up holding the same slice. Admissibility is the bare demand that you not be beaten in that comprehensive way. Notice how little has gone in: a scoring of outcomes state by state, and the refusal to be dominated. No weights over the states, no probabilities, have been mentioned.

On the other side stands the kind of rule we are trying to characterize, defined the textbook way using a probability distribution over the states. A probability distribution over the finite set $\Theta$ is just an assignment of a weight $\pi(\theta) \ge 0$ to each state with the weights summing to one, exactly the kind of object Section~\ref{sec:probabilities} forced into existence; here it represents how likely you take each state to be \emph{before} seeing your observation, which is traditionally called a \emph{prior}.

\begin{definition}[Bayes risk, Bayes-optimality, prior]
\label{def:bayes}
Given a prior $\pi$ over the states, the \emph{Bayes risk} of a rule $\delta$ is its prior-weighted average risk across states,
\[
r(\pi, \delta) \;=\; \sum_{\theta} \pi(\theta)\, R(\theta, \delta).
\]
A rule is \emph{Bayes-optimal} for $\pi$ if it makes this single number as small as possible. The prior is \emph{non-dogmatic} if it puts positive weight on every state, $\pi(\theta) > 0$ for all $\theta$, that is, it refuses to rule any state out in advance.
\end{definition}

A Bayes-optimal rule is exactly a maximizer of expected utility in disguise: it carries a probability distribution over the world ($\pi$), it scores outcomes (by the negative of the loss $L$), and it acts so as to make the probability-weighted score as good as possible. So the content of the coming theorem is that this familiar object, the expected utility maximizer, and the bare notion of being undominated, are the same thing.

\subsection{The theorem}

One direction is easy and worth seeing first, because it reveals why the non-dogmatic condition is the right one.

\begin{theorem}[Bayes-optimal rules are admissible]
\label{thm:cc-easy}
When $\Theta$ and $A$ are finite, any rule that is Bayes-optimal for a non-dogmatic prior $\pi$ is admissible.
\end{theorem}

\begin{proof}
Let $\delta$ minimize the Bayes risk $r(\pi, \cdot)$ for a prior $\pi$ that puts positive weight on every state, and suppose for contradiction that some $\delta^*$ Pareto-improves on $\delta$. Then $R(\theta, \delta) \ge R(\theta, \delta^*)$ for every $\theta$, with strict inequality for some particular state $\theta_0$. Multiply each of these comparisons by the corresponding weight $\pi(\theta)$ and add them up. Because every weight is positive, and the comparison at $\theta_0$ is strict with $\pi(\theta_0) > 0$, the sum comes out strictly larger on the $\delta$ side: $r(\pi, \delta) > r(\pi, \delta^*)$. That contradicts $\delta$ having the smallest Bayes risk. So no Pareto improvement exists and $\delta$ is admissible. The positivity of the prior is exactly what lets a strict improvement in the single state $\theta_0$ survive the averaging; a prior that assigned $\theta_0$ zero weight could shrug off an improvement that happens only there.
\end{proof}

The substantial direction is the converse, the complete class theorem proper, which says the Bayes-optimal rules are not merely \emph{some} of the admissible rules but \emph{all} of them: every way of being undominated is captured by some prior.

\begin{theorem}[Complete class theorem]
\label{thm:cc}
When $\Theta$ and $A$ are finite, a decision rule is admissible if and only if it is Bayes-optimal for some prior $\pi$ over the states.
\end{theorem}

Before the proof, the one piece of geometry it uses, stated for a reader who has not seen it. A set of points is \emph{convex} if, whenever it contains two points, it contains the entire straight segment joining them (a filled disk is convex; a crescent is not). The fact we need, the \emph{separating hyperplane theorem}, says that two convex sets that do not overlap can always be cleanly separated by a flat divider (a line in two dimensions, a plane in three, a ``hyperplane'' in general), with one set entirely on each side. That divider is described by the direction perpendicular to it, called its \emph{normal} direction, and the proof's whole trick is that this perpendicular direction, suitably normalized, is a prior.

\begin{proof}[Geometry of the proof]
Represent each available rule by its risk profile, the point $\big(R(\theta, \delta)\big)_{\theta}$ in the space with one axis per state. Allow the system to randomize, that is, to flip a private coin and follow one rule or another according to the result; randomizing between two rules produces the corresponding blend of their risk profiles, the point partway along the segment between them. So the set $\mathcal{R}$ of all achievable risk profiles, including these blends, is convex, and with finitely many basic rules it is also closed and bounded (it has a definite boundary and does not run off to infinity). Lower loss is better in every coordinate, so a system wants to sit as far toward the lower-left as the set allows, on its lower frontier.

Take an admissible rule $\delta$, sitting at the point $s = \big(R(\theta, \delta)\big)_\theta$. Admissibility says nothing in $\mathcal{R}$ beats $s$ in the strong sense: no achievable point lies at least as low as $s$ in every coordinate and strictly lower in some. Phrase this geometrically. The region of points that \emph{would} dominate $s$, the points sitting weakly below and to the left of $s$ in every coordinate, is a corner region hanging off $s$, and it is itself convex; admissibility says this dominating corner does not meet the achievable set $\mathcal{R}$. Two non-overlapping convex sets can be separated by a flat divider. The divider's perpendicular direction is a list of numbers $\pi = (\pi_\theta)_\theta$, one per state, and because the corner we separated away from extends toward smaller values in every coordinate, this perpendicular direction can be taken to point the same way in every coordinate, with all components nonnegative. Rescale it so its components sum to one, and it is precisely a prior over the states. What the separation asserts, read back in terms of that prior, is that $s$ achieves the smallest possible value of the prior-weighted sum $\sum_\theta \pi_\theta R(\theta, \cdot)$ over all of $\mathcal{R}$, which is to say $\delta$ is Bayes-optimal for $\pi$. Together with Theorem~\ref{thm:cc-easy}, admissibility and Bayes-optimality coincide.
\end{proof}

A two-state picture makes the geometry something you can see. When there are only two states, $\Theta = \{\theta_1, \theta_2\}$, each rule is a dot in an ordinary plane, its horizontal coordinate the risk in state one and its vertical coordinate the risk in state two. The dominated rules sit up and to the right, beaten by something below-and-left of them; the admissible rules form the lower-left boundary of the cloud of achievable points. A prior $(\pi_1, \pi_2)$ is a direction of tilt, and the Bayes-optimal rule for that prior is the boundary dot you reach by sliding a straight line of the corresponding slope as far down and to the left as it will go before leaving the cloud. As you swing the prior from ``all weight on state one'' to ``all weight on state two,'' the point of contact slides along the entire lower-left frontier, and every admissible dot is touched by exactly the prior whose line rests tangent there. The expected utility maximizers are not one virtuous corner of the space of good rules; traced out as the prior varies, they \emph{are} the frontier of undominated rules.

\subsection{What went in, and what came out}

The reason this is a \emph{foundation} and not merely a restatement is the accounting of inputs against outputs. Into the theorem went: finite lists of states and actions, a scoring of outcomes state by state (the loss $L$), and the single principle of admissibility, do not be dominated. Out of the theorem came, assumed nowhere in advance, a probability distribution over the states and the conclusion that the system's behavior maximizes expected utility weighted by that distribution. A genuine decision-theoretic foundation, in Demski's words, should not assume that the system has any probabilistic beliefs to begin with, and here the beliefs are not assumed; they are constructed, out of the geometry of which risk profiles are achievable. The probabilities and the utilities arrive together, the two readings of one separating divider, exactly as the intuitive arguments of Section~\ref{sec:probabilities} hinted but could not cleanly deliver.

Several things that look like assumptions can be removed, which reassures us about what is really essential. The likelihood $F$ is not truly needed: a noiseless observation is the special case where $F$ is always zero or one, and observation noise can always be folded into a finer description of the states (let a single ``state'' specify both the true situation and the reading on a faulty instrument). The randomizing used to make the achievable set convex need not be assumed to be probabilistic from outside; one can build an explicit coin into the description of the world and recover the same blends without presupposing the very notion of probability the theorem is meant to produce. And a technical axiom usually posited in this area, the so-called independence axiom of von Neumann and Morgenstern, here turns into a \emph{consequence}: any violation of it shows up as a Pareto improvement, so admissibility forces it rather than assuming it. The finiteness of the lists of states and actions is the assumption that does real work and is hardest to discard; with infinite possibilities the geometry becomes genuinely more delicate and the case for Bayesianism weakens, which Demski flags as an interesting frontier rather than a settled matter.

\subsection{A bonus: the same theorem argues for a kind of utilitarianism}

One pleasant surprise is that the complete class theorem reads twice. Reinterpret the index $\theta$ not as a possible state of one system's world but as a \emph{person} in a group, let the actions $A$ be the collective choices available to the group, and let $R(\theta, \delta)$ measure how badly off person $\theta$ ends up under collective policy $\delta$. Now admissibility becomes the requirement that the group not adopt a policy that another available policy beats for somebody while tying for everybody, which is the standard notion of social efficiency. And the prior $\pi$ the theorem manufactures becomes a set of weights, one per person. The theorem then says: every efficient social policy maximizes some positively-weighted sum of the individuals' (cardinal) welfare. This is essentially Harsanyi's aggregation theorem, an argument that a social planner who respects unanimous improvements must act like a weighted utilitarian, and it drops out of the same geometry for free. One caveat raised in discussion, due to Jessica Taylor, is that the theorem only certifies an efficient outcome as \emph{equal to} some weighted sum after the fact, rather than showing the outcome was \emph{produced by} computing that sum; the gap between ``describable as'' and ``mechanically does'' bites here exactly as it does in the single-agent reading, which is the subject of the final section.

\section{What these theorems do and do not say}
\label{sec:scope}

It is worth closing by stating precisely what has been established, because the coherence results are over-read about as often as they are used.

They describe behavior from the outside; they do not describe the machine inside. The complete class theorem says that any undominated system \emph{can be described} as maximizing expected utility for some prior. It does not say the system contains a prior, a utility function, or a little maximizing loop. A thermostat, a trained network, and an idealized planner might all be undominated in some decision problem and so all be summarizable by the same probabilities-and-utilities, while being built in utterly different ways inside. This is the precise content of the word ``behaves'' that we flagged at the very start. The theorems fix the shadow that effective goal-pursuit throws on behavior; they are silent about the object throwing it. That silence is the price of the robustness, since a description that holds whatever the internals are cannot, by the same token, tell you what the internals are. So coherence arguments by themselves cannot tell you how to build a safe system, nor how to read a system's goals off its parts. What they can tell you is the target that outward behavior drifts toward as a system gets better at pursuing its goals, and therefore what a more capable system will tend to look like from where we stand: more like a clean expected utility maximizer, with fewer drainable inconsistencies, as the pressure to be effective rises.

They lean on a ruler we cannot yet find in general. Every notion of being dominated, including the loss function the complete class theorem treats as given, presupposes a currency in which being dominated is bad. Section~\ref{sec:measuringstick} is the standing reminder that for an arbitrary physical object we do not know how to identify that currency without quietly assuming the very agency we were hoping to detect. Solving the measuring-stick problem, finding a world's resources from its physics alone, would promote these theorems from a description we apply once we already know what counts as a resource into a test we could run on raw physics to locate the goal-pursuers in it. That is squarely the business of descriptive agent foundations, and the handoff is clean: the coherence theorems supply the target shape that effective agents present from outside, and the descriptive and physical programs ask what real pressures produce that shape and what physical quantity actually plays the role of money.

They assume an idealization that real systems break, and the breakages are where the subject goes next. Perfect coherence cannot be computed, so every real system is incoherent somewhere; a system that is uncertain about the consequences of its own beliefs (that knows the rules of arithmetic, say, but has not worked out whether some particular large number is prime) cannot even assign the sharp probabilities the Dutch book demands, which is the precise crack that the theory of logical uncertainty is built to fill. And the tidy assumptions behind the cleanest theorems, that preferences are complete, fixed, and independent of how you arrived at the current choice, fail for humans and for trained systems alike, which is exactly the gap that more realistic, more mechanistic results are meant to close. The coherence theorems are best held not as the last word on what an agent is but as its sturdiest first foothold: the part of effective goal-pursuit that is forced, structurally, stated cleanly enough that we can see precisely which of its assumptions to attack next.

\section*{Sources}
\addcontentsline{toc}{section}{Sources}

\begin{itemize}[leftmargin=1.5em]
\item Eliezer Yudkowsky, \emph{Coherent decisions imply consistent utilities} (sections read: Introduction; Why not circular preferences?; Probabilities and expected utilities through the Conditional probability subsection; Conclusion). \url{https://www.lesswrong.com/s/SgomvxZ3cJWy2SBCu/p/RQpNHSiWaXTvDxt6R}
\item John Wentworth, \emph{The measuring stick of utility problem}. \url{https://www.lesswrong.com/posts/73pTioGZKNcfQmvGF/the-measuring-stick-of-utility-problem}
\item Abram Demski, \emph{Complete Class: Consequentialist Foundations}. \url{https://www.lesswrong.com/posts/sZuw6SGfmZHvcAAEP/complete-class-consequentialist-foundations}
\end{itemize}

\end{document}
